\documentclass[11pt]{article}
\usepackage[margin=28mm]{geometry}
\usepackage{amsmath,amssymb,amsthm,mathtools,mathrsfs,booktabs,longtable,array}
\usepackage[colorlinks]{hyperref}
\newcommand{\statement}[3]{\subsection*{#1. #3 [#2]}}
\begin{document}
\statement{C.30}{Proved}{Unit hexagon solutions are full automorphisms}
Let $F=\widehat F(x,y)$, put $z=(xy)^{-1}$, and let
$\lambda\in\widehat{\mathbb Z}^{\times}$ and $f\in F'$.
Write $m=(\lambda-1)/2\in\widehat{\mathbb Z}$.
Suppose the following are identities in the full profinite group:
\[
 \begin{split}
 f(x,y)f(y,x)&=1,\\
 f(x,y)x^m f(z,x)z^m f(y,z)y^m&=1.
 \end{split}                                                    \tag{C.30.1}
\]
Then the continuous endomorphism
\[
             \varphi(x)=x^\lambda,\qquad
             \varphi(y)=f^{-1}y^\lambda f
\]
is an automorphism. In particular no separate invertibility hypothesis
is needed for a unit pair satisfying these two full equations. If it
also satisfies the full pentagon, it belongs to the full
Grothendieck--Teichm\"uller group, including its unit requirement.

\begin{proof}
This is an algebraic consequence of C.29, but it is essential not to
assume invertibility while deriving the third-cusp equation.
Set $A=f(x,y)$, $B=f(y,z)$ and $C=f(z,x)$. The second equation in
(C.30.1) first gives
\[
                C z^m B=x^{-m}A^{-1}y^{-m}.                      \tag{C.30.2}
\]
Apply that same identity under the continuous substitution
$(x,y)\mapsto(x,z)$. Its third argument is
\[
                w=(xz)^{-1}=xyx^{-1}=z^{-1}yz.
\]
The first equation in (C.30.1), applied under substitutions, and
conjugation-equivariance of profinite words give
\[
 f(x,z)=C^{-1},\qquad
 f(w,x)=xA^{-1}x^{-1},\qquad
 f(z,w)=z^{-1}B^{-1}z.
\]
Consequently the substituted hexagon becomes
\[
 C^{-1}x^{m+1}A^{-1}y^{m+1}B^{-1}z^{m+1}=1,
\]
where $x^{-1}z^{-1}=y$ was used. Rearrangement gives the second identity
\[
          C z^{-m-1}B=x^{m+1}A^{-1}y^{m+1}.                     \tag{C.30.3}
\]
Multiply (C.30.3) by the inverse of (C.30.2). Since $2m+1=\lambda$,
then conjugate the resulting equality by $x^m$. One obtains
\[
 x^\lambda A^{-1}y^\lambda A
       =(x^mC)(xy)^\lambda(x^mC)^{-1}.                         \tag{C.30.4}
\]
All substitutions used here are legitimate continuous homomorphisms
out of the free profinite group; none has been assumed invertible.
Equation (C.30.4) is the third peripheral condition, and C.29 makes
$\varphi$ a full automorphism.

For completeness the pentagon assertion concerns actual units, not
only endomorphisms. The normalized two-equation pairs are now
contained in $\operatorname{Aut}(F)$. Their pentagon solutions form
a closed submonoid there, by the parenthesized-braid presentation
\cite[Lemma 7.1 and Proposition 7.3]{CEBHR2019}. The automorphism group
of a finitely generated profinite group is profinite: characteristic
open subgroups are cofinal, and their finite automorphism groups
give its topology. Thus any closed submonoid is a group, since
$\varphi^{n!-1}\to\varphi^{-1}$. This supplies the inverse within
the same pentagon solution set.
\end{proof}

\paragraph{Source discrepancy and scope.}
Lochak and Schneps explicitly asked in 2006 whether a unit normalized
pair is automatically invertible, and then asked the same question
with their relations (I), (II), (III) imposed
\cite[Question 1.1]{A_LS2006}. Fresse's later survey also distinguishes
the profinite invertibility issue from the Malcev case
\cite[\S1.1, manuscript p.~14]{CEFresse2019}.
C.30 answers the constrained version by the displayed full-group
calculation and C.29's finite-cover proof. It does not use the
pro-$p$ Frattini criterion, infer invertibility merely from
abelianization, or claim a literature priority. It also does not
prove that a bare profinite pentagon has a unit parameter or implies
the hexagon; those are different requirements in A.37--A.39.

The unconstrained version of that old question has a small finite
counterexample, which also calibrates this source comparison.
Take $\lambda=1$, $f=[x,y]^2$ and the quotient onto $A_5$ with
$x=(123)$ and $y=(345)$, using $[x,y]=xyx^{-1}y^{-1}$ and
right-to-left permutation products. Then $f$ has image $(134)$,
and $f^{-1}yf$ has image $(135)$. The original pair generates
$A_5$, whereas the substituted pair generates the order-twelve
point stabilizer fixing $4$. Thus this unrestricted normalized
endomorphism is not surjective. The exact finite calculation is
recorded in \path{code/check_unconstrained_substitution.py}.

\end{document}
